\graphicspath{ {./lab1/images/part1/} }

\newpage
\section{Лабораторная работа № 1, часть 1: Генерация псевдослучайных чисел}
\subsection{Инициализация ГПСЧ для X $\sim$ U[0, 1]}

Инициализируем генератор \texttt{rng}, используя конструктор \texttt{default\_rng} из библиотеки \texttt{numpy}:
\begin{minted}{python}
rng = np.random.default_rng()
\end{minted}



\subsection{Нахождение зависимости между мат ожиданием, дисперсией и размером выборки}

Для исследования сходимости выборочных характеристик генерируются массивы различной длины и вычисляются их математические ожидания и дисперсии.

Построим график (Рис. \ref{fig:mean_vs_N}) зависимости математического ожидания \texttt{mean} от размера выборки $N$.
\begin{figure}[htbp]
    \centering
    \includegraphics[width=0.6\textwidth]{mean_vs_N.png}
    \caption{Зависимость выборочного математического ожидания от размера выборки $N$.}
    \label{fig:mean_vs_N}
\end{figure}

Построим график (Рис. \ref{fig:var_vs_N}) зависимости дисперсии \texttt{var} от размера выборки $N$.
\begin{figure}[htbp]
    \centering
    \includegraphics[width=0.6\textwidth]{var_vs_N.png}
    \caption{Зависимость выборочной дисперсии от размера выборки $N$.}
    \label{fig:var_vs_N}
\end{figure}



\newpage
\subsection{Построение гистограммы и вычисление погрешности (расхождение сгенерированных чисел с теоретическим значением)}

Генерируется выборка большого объема ($N = 10^8$) и строится ее нормированная гистограмма:
\begin{minted}{python}
sample_size = 10 ** 8
sample = rng.uniform(0., 1., size=sample_size)
num_bins = 32
sample_hist, bin_edges = np.histogram(sample, bins=num_bins, density=True)
\end{minted}

Построим приближенную гистограмму (Рис. \ref{fig:histogram}) равномерного распределения на основе сгенерированных данных (\texttt{sample}).
\begin{figure}[htbp]
    \centering
    \includegraphics[width=0.6\textwidth]{histogram.png}
    \caption{Приближенная гистограмма равномерного распределения $U[0, 1]$.}
    \label{fig:histogram}
\end{figure}

Определим функцию для вычисления погрешности приближения гистограммы к теоретической плотности ($p_{\text{true}} = 1$):
\begin{minted}{python}
def distribution_error(p: np.ndarray, p_true: float = 1.):
    return np.sum((p - p_true) ** 2)
\end{minted}

Построим график (Рис. \ref{fig:error_vs_N}) зависимости погрешности \texttt{distribution\_error} от размера выборки $N$.
\begin{figure}[htbp]
    \centering
    \includegraphics[width=0.6\textwidth]{error_vs_N.png}
    \caption{Зависимость погрешности распределения от размера выборки $N$.}
    \label{fig:error_vs_N}
\end{figure}



\newpage
\subsection{Вычисление коэффициента корреляции между двумя выборками из одного распределения, но построенными разными способами}

Для оценки линейной зависимости реализуем вычисление коэффициента корреляции Пирсона:
\begin{minted}{python}
def pearson_r(x: np.ndarray, y: np.ndarray):
    return np.sum((x - x.mean()) * (y - y.mean())) / np.sqrt(np.sum((x - x.mean()) ** 2) * np.sum((y - y.mean()) ** 2))
\end{minted}

\vspace{0.5em}
\textbf{(а)} генерация двух независимых выборок.

Построим график (Рис. \ref{fig:corr_a}) зависимости коэффициента корреляции \texttt{pearson\_r} от размера выборки $N$ (вариант а).
\begin{figure}[htbp]
    \centering
    \includegraphics[width=0.6\textwidth]{corr_a.png}
    \caption{Зависимость коэффициента корреляции от размера выборки $N$ (вариант а).}
    \label{fig:corr_a}
\end{figure}

\vspace{0.5em}
\textbf{(б)} генерация единой последовательности и разделение на четные/нечетные элементы.

Построим график (Рис. \ref{fig:corr_b}) зависимости коэффициента корреляции \texttt{pearson\_r} от размера выборки $N$ (вариант б).
\begin{figure}[htbp]
    \centering
    \includegraphics[width=0.6\textwidth]{corr_b.png}
    \caption{Зависимость коэффициента корреляции от размера выборки $N$ (вариант б).}
    \label{fig:corr_b}
\end{figure}
